\cleardoublepage
\chapter{Pelaksanaan Pekerjaan Bulan Agustus}
\label{chap:pekerjaan-agustus}
\section{Gambaran Umum Pekerjaan}
Pada Agustus 2026, kegiatan kerja praktik lebih banyak diarahkan pada dukungan eksplorasi solusi, pengumpulan informasi teknis-komersial, dan dokumentasi hasil pembahasan bersama mitra. Fokus pekerjaan yang paling jelas untuk diuraikan sebagai proses adalah eksplorasi solusi Penerangan Jalan Umum Tenaga Surya (PJUTS) melalui materi produk dan diskusi dengan PT Jarwinn Solar Panel.

\section{Proses Eksplorasi Solusi PJUTS}
\subsection{Pengenalan Alternatif Solusi PJUTS}
Pada kegiatan ini, penulis mengikuti diskusi bersama mentor dan Om Jun selaku Manager Product Tech, serta perwakilan PT Jarwinn Solar Panel yang membahas alternatif solusi Penerangan Jalan Umum Tenaga Surya (PJUTS). Materi yang dibahas meliputi jenis konfigurasi sistem tenaga surya, yaitu \textit{on-grid}, \textit{off-grid}, dan \textit{hybrid}, serta variasi produk lampu PJUTS dengan tipe \textit{2-in-1} dan \textit{all-in-one}. Melalui diskusi ini, penulis memperoleh gambaran awal mengenai perbedaan konfigurasi dan karakteristik produk yang dapat dipertimbangkan sesuai dengan kebutuhan proyek.

\subsection{Dokumentasi Hasil Diskusi}
Setelah mengikuti diskusi mengenai alternatif solusi PJUTS, penulis menyusun dokumen \textit{Minutes of Meeting} (MoM) untuk merangkum informasi dan poin-poin pembahasan yang disampaikan oleh pihak terkait. Selain itu, hasil diskusi juga dirangkum ke dalam satu slide presentasi yang memuat informasi utama sebagai bahan dokumentasi internal.

\section{Observasi Teknologi pada Pameran DTI}
\subsection{Pelaksanaan Observasi}
Penulis dan mentor menghadiri kunjungan untuk mewakili Telkom Group dalam acara pameran \textit{Digital Transformation Indonesia}. Selama kegiatan, terdapat berbagai teknologi dan solusi inovatif yang ditampilkan oleh masing-masing vendor. Penulis juga mendokumentasikan kunjungan sebagai bahan referensi. Aktivitas ini memberikan pengetahuan mengenai pendekatan produk di lingkungan industri.

\subsection{Hasil Dokumentasi}
Hasil kegiatan dirangkum dalam bentuk dokumen catatan dan dokumentasi foto. Penulis tidak melakukan observasi teknis terhadap produk yang ditampilkan, sehingga hasil pengamatan tersebut digunakan sebagai bahan eksplorasi sekaligus referensi.

\section{Penyusunan dan Penyampaian Presentasi Akhir}
\subsection{Menyusun Materi Pekerjaan}
Menjelang akhir masa kerja praktik, penulis menyiapkan bahan untuk menyampaikan pekerjaan yang telah dijalani, terutama eksplorasi \textit{Smart Depot} berupa perancangan \textit{Operational Dashboard} pada bulan Juli serta kegiatan eksplorasi solusi pada bulan Agustus. Penyusunan materi diarahkan oleh mentor, kemudian didokumentasikan dalam bentuk PowerPoint dan lembar kerja berupa \textit{log activity} yang disertai laporan singkat untuk perusahaan.

\subsection{Presentasi}
Untuk mengakhiri kegiatan kerja praktik, penulis ditugaskan untuk menyampaikan presentasi akhir kepada mentor dan beberapa staf BNPD, termasuk tim \textit{Research \& Development} (R\&D) dan \textit{Technology Solution}. Presentasi ini bertujuan untuk memaparkan hasil pekerjaan yang telah dilakukan selama dua bulan kerja praktik. Materi utama yang disampaikan mencakup eksplorasi konsep \textit{Smart Depot Management System} (SDMS), perkembangan perancangan \textit{Operational Dashboard}, serta gambaran manfaat penerapan solusi tersebut dalam mendukung operasional depo. Penulis juga memberikan penjelasan singkat mengenai hasil diskusi alternatif Penerangan Jalan Umum Tenaga Surya (PJUTS) bersama PT Jarwinn Solar Panel.

\section{Tantangan dan Penyelesaian Masalah}
Sejumlah tantangan dialami penulis saat melaksanakan kerja praktik pada bulan Agustus. Penulis masih perlu mempelajari berbagai istilah teknis dan produk yang dibahas karena belum familiar, khususnya Penerangan Jalan Umum Tenaga Surya (PJUTS) yang dibahas bersama PT Jarwinn Solar Panel. Banyaknya informasi yang disampaikan selama diskusi membuat penulis harus mempelajari kembali perbedaan karakteristik produk, pilihan konfigurasi, dan kesesuaiannya dengan kebutuhan proyek. Tantangan ini juga dirasakan ketika menyusun dokumen \textit{Minutes of Meeting} (MoM) dan satu slide presentasi, karena poin-poin pembahasan perlu dirangkum secara jelas sesuai dengan standar perusahaan agar mudah dipahami. Untuk mengatasinya, penulis meninjau kembali catatan diskusi dan melakukan konfirmasi kepada mentor.
